% =====================================================================
%  1bat-ud3-12.tex  —  HORA 12: Gimnàstica d'anàlisi de funcions (mecanització)
%  (sense preàmbul: s'inclou des de main.tex)
% =====================================================================
\horatitol{Sessió 12 — Gimnàstica d'anàlisi de funcions (mecanització)}%
{Treball simbòlic per automatitzar dos procediments: el domini teòric i els punts de tall amb els eixos.}

\subsection{Idea clau}

Després de donar sentit a les funcions en contextos reals, avui toca
\textbf{guanyar destresa}. Treballarem de manera \emph{simbòlica} (sense
històries): es tracta que el càlcul del \textbf{domini} i dels \textbf{punts de
tall} surti automàtic. És com fer escales abans de tocar una peça.

\begin{idea}
\textbf{Domini teòric: quins valors de $x$ \emph{no} es poden posar?}
\begin{itemize}[leftmargin=1.4em, itemsep=2pt]
  \item \textbf{Polinomis} (rectes, paràboles, $x^3$\dots): es pot posar
        \textbf{qualsevol} $x$. Domini $=\mathbb{R}$.
  \item \textbf{Fraccions} $\dfrac{P(x)}{Q(x)}$: el denominador \textbf{no pot ser
        zero}. S'exclouen els valors que anul·len $Q(x)$.
  \item \textbf{Arrels quadrades} $\sqrt{P(x)}$: el que hi ha dins \textbf{no pot ser
        negatiu}. Cal $P(x)\ge 0$.
\end{itemize}
\textbf{Punts de tall:}
\begin{itemize}[leftmargin=1.4em, itemsep=2pt]
  \item Amb l'eix $Y$: fes $x=0$ (si $0$ és al domini), punt $(0,f(0))$.
  \item Amb l'eix $X$: fes $f(x)=0$ i resol.
\end{itemize}
\textbf{Paranys de sempre:} oblidar les restriccions del domini i confondre quin eix
dóna cada tall.
\end{idea}

\subsection{Exemples}

\begin{exemple}[domini d'una fracció]
Sigui $f(x)=\dfrac{x+1}{x-3}$. El denominador s'anul·la quan $x-3=0$, és a dir
$x=3$. Per tant, \textbf{tots els valors menys el $3$}:
\[
\mathrm{Dom}\,f=\mathbb{R}\setminus\{3\}.
\]
El tall amb l'eix $Y$ ($x=0$): $f(0)=\dfrac{1}{-3}=-\dfrac{1}{3}$, punt
$\left(0,-\tfrac{1}{3}\right)$. El tall amb l'eix $X$ ($f(x)=0$): una fracció és
zero quan el \emph{numerador} ho és, $x+1=0 \Rightarrow x=-1$, punt $(-1,0)$.
\end{exemple}

\begin{exemple}[domini d'una arrel]
Sigui $g(x)=\sqrt{x-2}$. El que hi ha dins l'arrel ha de ser $\ge 0$:
$x-2\ge 0 \Rightarrow x\ge 2$. Per tant
\[
\mathrm{Dom}\,g=[2,+\infty).
\]
No hi ha tall amb l'eix $Y$ ($x=0$ no és al domini). Tall amb l'eix $X$:
$\sqrt{x-2}=0 \Rightarrow x-2=0 \Rightarrow x=2$, punt $(2,0)$.
\end{exemple}

\subsection{Exercicis resolts}

\begin{exresolt}[anàlisi completa d'una paràbola]
Per a $f(x)=x^2-x-6$, troba el domini, el tall amb l'eix $Y$ i els talls amb l'eix
$X$.
\solucio
\textbf{Domini:} és un polinomi, així que $\mathrm{Dom}\,f=\mathbb{R}$.

\textbf{Tall amb $Y$} ($x=0$): $f(0)=-6$, punt $(0,-6)$.

\textbf{Talls amb $X$} ($f(x)=0$): $x^2-x-6=0$. Factoritzem buscant dos nombres que
multiplicats donin $-6$ i sumats $-1$: són $-3$ i $2$. Així $(x-3)(x+2)=0$, d'on
$x=3$ i $x=-2$. Punts $(3,0)$ i $(-2,0)$.
\resposta{$\mathrm{Dom}=\mathbb{R}$; tall $Y$: $(0,-6)$; talls $X$: $(-2,0)$ i $(3,0)$.}
\end{exresolt}

\begin{exresolt}[anàlisi d'una funció racional]
Per a $f(x)=\dfrac{2x-4}{x+1}$, troba el domini i els punts de tall amb els dos eixos.
\solucio
\textbf{Domini:} el denominador s'anul·la si $x+1=0 \Rightarrow x=-1$. Per tant
$\mathrm{Dom}\,f=\mathbb{R}\setminus\{-1\}$.

\textbf{Tall amb $Y$} ($x=0$): $f(0)=\dfrac{-4}{1}=-4$, punt $(0,-4)$.

\textbf{Tall amb $X$} ($f(x)=0$): la fracció és zero quan el numerador ho és:
$2x-4=0 \Rightarrow x=2$ (i $2$ sí que és al domini). Punt $(2,0)$.
\resposta{$\mathrm{Dom}=\mathbb{R}\setminus\{-1\}$; tall $Y$: $(0,-4)$; tall $X$: $(2,0)$.}
\end{exresolt}

\subsection{Exercicis per fer}

\noindent\textit{Bateria graduada. Per a cada funció: domini, tall amb l'eix $Y$ i
talls amb l'eix $X$. Treballa sense calculadora sempre que puguis.}

\begin{exercicis}
\begin{exlist}
  \item \textbf{Nivell bàsic (rectes i paràboles).}
    \begin{enumerate}[label=\alph*), leftmargin=1.6em, topsep=2pt]
      \item $f(x)=2x-6$
      \item $f(x)=x^2-9$
      \item $f(x)=x^2-5x+6$
    \end{enumerate}
  \item \textbf{Nivell mitjà (més paràboles).}
    \begin{enumerate}[label=\alph*), leftmargin=1.6em, topsep=2pt]
      \item $f(x)=-x^2+4x$
      \item $f(x)=x^2+2x-8$
      \item $f(x)=2x^2-8$
    \end{enumerate}
  \item \textbf{Nivell avançat (fraccions).} Recorda: denominador $\neq 0$.
    \begin{enumerate}[label=\alph*), leftmargin=1.6em, topsep=2pt]
      \item $f(x)=\dfrac{x-5}{x-2}$
      \item $f(x)=\dfrac{3x}{x+4}$
    \end{enumerate}
  \item \textbf{Nivell avançat (arrels).} Recorda: el radicand $\ge 0$.
    \begin{enumerate}[label=\alph*), leftmargin=1.6em, topsep=2pt]
      \item $f(x)=\sqrt{x+3}$
      \item $f(x)=\sqrt{2x-6}$
    \end{enumerate}
  \item \textbf{(Parany)} Un company diu que $f(x)=\dfrac{x+2}{x^2-4}$ té domini
        $\mathbb{R}\setminus\{4\}$. Troba el domini correcte. (Pista: factoritza el
        denominador.)
  \item \textbf{(Correcció creuada)} Intercanvia la teva bateria amb la d'una altra
        parella i corregiu-vos mútuament. Anoteu els errors que trobeu i classifiqueu-los:
        són de \emph{domini} (oblidar una restricció) o de \emph{talls} (confondre els
        eixos o equivocar-se en resoldre l'equació)?
\end{exlist}
\end{exercicis}

% ---------------------------------------------------------------------
\fullsolucions{Sessió 12}
\begin{solucions}
\begin{sollist}
  \item (a) $\mathrm{Dom}=\mathbb{R}$; $Y$: $(0,-6)$; $X$: $(3,0)$. \quad
        (b) $\mathrm{Dom}=\mathbb{R}$; $Y$: $(0,-9)$; $X$: $(3,0)$ i $(-3,0)$. \quad
        (c) $\mathrm{Dom}=\mathbb{R}$; $Y$: $(0,6)$; $X$: $(2,0)$ i $(3,0)$.
  \item (a) $\mathrm{Dom}=\mathbb{R}$; $Y$: $(0,0)$; $X$: $(0,0)$ i $(4,0)$. \quad
        (b) $\mathrm{Dom}=\mathbb{R}$; $Y$: $(0,-8)$; $X$: $(-4,0)$ i $(2,0)$. \quad
        (c) $\mathrm{Dom}=\mathbb{R}$; $Y$: $(0,-8)$; $X$: $(2,0)$ i $(-2,0)$.
  \item (a) $\mathrm{Dom}=\mathbb{R}\setminus\{2\}$; $Y$: $\left(0,\tfrac{5}{2}\right)$;
        $X$: $(5,0)$. \quad
        (b) $\mathrm{Dom}=\mathbb{R}\setminus\{-4\}$; $Y$: $(0,0)$; $X$: $(0,0)$.
  \item (a) $\mathrm{Dom}=[-3,+\infty)$; $Y$: $(0,\sqrt{3})$; $X$: $(-3,0)$. \quad
        (b) $\mathrm{Dom}=[3,+\infty)$; $Y$: no en té ($0\notin\mathrm{Dom}$);
        $X$: $(3,0)$.
  \item El denominador és $x^2-4=(x-2)(x+2)$, que s'anul·la a $x=2$ i $x=-2$. Per
        tant $\mathrm{Dom}=\mathbb{R}\setminus\{-2,2\}$ (no $\{4\}$).
  \item Resposta oberta (correcció entre parelles). Es valora classificar bé cada
        error com a error de domini o de talls.
\end{sollist}
\end{solucions}
